% BEGIN R1_ACTION_NATIVE_GRADIENT_DISPERSION_V1
\subsection{Action-native gradient and dispersion audit}
\label{sec:action-native-gradient-dispersion}

The remaining linear/quadratic stability question was tested directly in the
same constraint-reduced ADM plus Schutz--Sorkin action used for the
matter/radiation evolution. No observational amplitude, background parameter,
local-source normalization, or Gate-7 quantity enters this audit. For scalar
harmonic number $n\ge2$ we use the closed-$S^3$ eigenvalue
$q_n=n(n+2)$ and continue only the harmonic factors of the already derived
quadratic action. At $n=2$ this construction reproduces both the unreduced
action blocks and the reduced $(M,L,V)$ blocks with zero difference to
floating precision.

Writing the reduced equations as
\begin{equation}
M\ddot x+\Gamma\dot x+\Omega x=0,
\end{equation}
the instantaneous dispersion pencil is
\begin{equation}
\det\!\left[s^2M+s\Gamma+\Omega\right]=0,
\qquad
\Gamma=3HM+\dot M+L-L^T,
\qquad
\Omega=3HL+\dot L-V .
\end{equation}
This formulation retains the time dependence of the background and the
antisymmetric velocity mixing rather than inferring stability from the kinetic
matrix alone.

The reduced kinetic matrix remains positive throughout the explicit scan.
The original audit through $n=320$ has minimum eigenvalue
$1.94099322\times10^{-7}$. An additional promotion check on the extended
principal-symbol sequence through $n=1280$ also remains
positive, with minimum eigenvalue $1.21879873\times10^{-8}$ at
$z=0.000$, $n=1280$.

The high-$k$ diagnostic separates ordinary pressureless-matter growth from a
gradient instability by testing the scaled quantity
$\mathrm{Re}(s)/\sqrt{k_{\rm phys}^2}$. Extending the harmonic sequence
from $(40,80,160,320)$ to $(80,160,320,640,1280)$ reduces the extrapolated
positive intercept from $3.79858666\times10^{-3}$ to $3.84535287\times10^{-4}$,
supporting convergence toward zero rather than a persistent $O(k)$
instability. In the extended run the non-radiation propagating branch has
\begin{equation}
c_{\rm scalar}^2 \ge 0.910373,
\end{equation}
while the radiation branch recovers $c_r^2=1/3$ with maximum absolute
extrapolation error $3.34943277\times10^{-4}$ over $0\le z\le2.1$.

As an independent closed-Horndeski cross-check, the pure gravity/scalar
Akama--Kobayashi coefficients remain positive:
\begin{equation}
G_{S,2}^{\rm min}=1.47147014\times10^{-3},
\qquad
F_{S,2}^{\rm min}=5.03726732\times10^{-1},
\qquad
\min(F_{S,2}/G_{S,2})
=2.590796.
\end{equation}

We therefore find no ghost or high-$k$ scalar gradient instability in the
audited frozen R1 domain at linear/quadratic order. This is a numerical
principal-symbol stability result for the effective closed-FLRW realization;
it is not a nonlinear-fluid theorem, a proof of ultraviolet completion, or a
microscopic completion of BHSM.
% END R1_ACTION_NATIVE_GRADIENT_DISPERSION_V1
